\section{问题分析}
\subsection{预测、购电计划与实时执行的分工}
问题一的输入曲线已知，可以直接优化全天购电和储能安排。问题二至四需要先预测再规划，实际需求偏离预测时，还要调整充放电并补购缺少的电力。本文据此安排\emph{预测、常规购电规划和实时执行}三个环节。一次规划的最优解只对应当次预测，全年效果还取决于预测误差及后续运行。由于储能可以在时段间转移电量，同样的误差在不同库存和电价下会产生不同支出，故主要按实际费用评价策略，并以预测精度辅助解释。

\subsection{四个问题的递进关系}
问题一先建立电量平衡、储能损耗及日循环约束，再判断充放电互斥条件能否作等价线性松弛。问题二沿用这些物理关系，根据历史预测误差修正净负荷，以应对紧急购电价格较高的风险，并在运行中根据实际供需调节储能。

问题三既增加光伏预报，又允许日内修改购电，需要通过对照辨别两者各自的作用。多次修改的费用还取决于合约如何累计调整量，因而需比较最终净调整和逐笔成交两种解释。问题四进一步预测未来电价，使购电和储能安排能随预期价格变化，并按交付时的实际电价检验费用。

\subsection{方法选择及相关研究}

如何在不同时段分配有限储电量，是光伏储能经济调度及日前、日内协调优化的共同问题\upcite{wu2018tou,xiao2019mpc}。模型预测控制在获得新观测后重新预测并优化未来一段时间，每次只执行第一个动作。确定性滚动控制采用一组预测，多情景鲁棒控制则考虑多种可能情况\upcite{wytock2017dem}。动态电价研究还将实际可执行的预测控制，与已知未来全部数据时的表现作对照\upcite{perez2023hem}。

多时间尺度模型预测控制还可结合分布鲁棒优化，处理未来场景及其概率分布的不确定性\upcite{li2024mpc}。本文根据历史误差分位修正各段净负荷，再用线性规划制定购电计划，计算规模较小。针对这种近似可能带来的影响，分别比较历史预测、不同预报来源、交易合约及储能反馈方式。经验分位的含义和限制在问题二中具体说明。
